\documentclass[11pt,a4paper]{coverletter}

\senderblock
  {Radheem Bin Razi}
  {Distributed Systems \& Cloud-Native Engineer}
  {Ilmenau, Germany \textbar\ \href{mailto:sheikh.radheem@gmail.com}{sheikh.radheem@gmail.com} \textbar\ \href{https://radheem.github.io/my-cv}{portfolio}}

\begin{document}

%% ===== ENGLISH =====
\recipient{Verivox}{Hiring Team}{}

\opening{Dear Hiring Team,}

\vspace{8pt}\noindent\textbf{1. Why Verivox?}\par\vspace{3pt}

A full-stack role across a high-traffic comparison platform, with room to own frontend and backend work end-to-end, and to shape AI-powered features that make better recommendations possible. For an engineer who takes pride in code quality, performance, and real user outcomes, this is a great opportunity. Verivox's mission to make markets transparent is something I personally find worth working for, since it makes a difference in the lives of millions of everyday people — I believe technology at its best does exactly this.

\vspace{8pt}\noindent\textbf{2. Why me?}\par\vspace{3pt}

\bullets

  \item Delivered features end-to-end using TypeScript across the stack, building React/Next.js frontends and NestJS/Node.js backends that connect to PostgreSQL and MongoDB through well-designed REST APIs (SeedLabs).

  \item Contributed to a high-throughput trading platform frontend and backend, driving performance, reliability and monitoring with DataDog and Grafana, while taking end-to-end ownership through CI/CD-style release coordination (Bluefin Exchange).

  \item Maintained a product-minded focus on performance optimization and code quality, consistently contributing to system-design discussions and rigorous code reviews to ensure scalable, maintainable architecture.

  \item Engineered AI-powered features by developing an LLM-driven automation suite (\texttt{cv-tailor}) with a FastMCP server, data ingestion pipelines, and bidirectional Google Sheets sync, directly aligning with your goal to integrate AI capabilities into consumer platforms.

  \item Architected and deployed cloud-native solutions by Dockerizing microservices and managing Kubernetes clusters with PostgreSQL, ensuring robust infrastructure aligned with modern DevOps practices.

\bulletsend

While my recent frontend work has primarily utilized React and Next.js, both are TypeScript-based component frameworks, and I am fully comfortable and eager to adopt Angular, RxJS, and Reactive Forms. My extensive NestJS and Node.js backend experience aligns directly with your stack, and I thrive in cross-functional Agile environments where I can collaborate closely with UX designers, QA engineers, and product owners to translate Figma designs into production-ready code.

I am available full-time immediately and open to part-time arrangements, based in Germany with relocation readiness within two to three weeks, and hold Blue Card eligibility requiring only a standard contract issuance without sponsorship overhead.

\closing{Sincerely,}{Radheem Bin Razi}

\clearpage
\selectlanguage{ngerman}

%% ===== DEUTSCH =====
\recipient{Verivox}{Personalabteilung}{}

\opening{Sehr geehrtes Hiring-Team,}

\vspace{8pt}\noindent\textbf{1. Warum Verivox?}\par\vspace{3pt}

Eine Full-Stack-Position auf einer stark frequentierten Vergleichsplattform, die die Möglichkeit bietet, Frontend- und Backend-Entwicklung von Anfang bis Ende zu übernehmen und KI-gestützte Funktionen zu gestalten, die präzisere Empfehlungen ermöglichen. Für eine Entwicklerin/einen Entwickler, der Wert auf Code-Qualität, Performance und messbare Nutzererfolge legt, ist dies eine hervorragende Chance. Die Mission von Verivox, Märkte transparent zu machen, ist ein Ziel, wofür ich mich gerne engagiere, da es das Leben von Millionen von Verbraucherinnen und Verbrauchern konkret verbessert – ich bin überzeugt, dass Technologie genau das in ihrer besten Form leistet.

\vspace{8pt}\noindent\textbf{2. Warum ich?}\par\vspace{3pt}

\bullets

  \item Umsetzung von Features von Anfang bis Ende mit TypeScript über den gesamten Stack hinweg, Entwicklung von React/Next.js-Frontends und NestJS/Node.js-Backends, die über gut strukturierte REST-APIs mit PostgreSQL und MongoDB verbunden sind (SeedLabs).

  \item Mitwirkung an Frontend und Backend einer hochperformanten Handelsplattform, Optimierung von Performance, Zuverlässigkeit und Monitoring mittels DataDog und Grafana sowie Übernahme der End-to-End-Verantwortung durch CI/CD-ähnliche Release-Koordination (Bluefin Exchange).

  \item Produktorientierter Fokus auf Performance-Optimierung und Code-Qualität, kontinuierliche Mitwirkung an Systemdesign-Diskussionen und strengen Code Reviews zur Sicherstellung einer skalierbaren und wartbaren Architektur.

  \item Entwicklung von KI-gestützten Funktionen durch Erstellung einer LLM-basierten Automatisierungs-Suite (\texttt{cv-tailor}) mit FastMCP-Server, Data-Ingestion-Pipelines und bidirektionalem Google-Sheets-Sync, was direkt Ihrem Ziel entspricht, KI-Fähigkeiten in Consumer-Plattformen zu integrieren.

  \item Architektur und Deployment cloudbasierter Lösungen durch Dockerisierung von Microservices und Management von Kubernetes-Clustern mit PostgreSQL, Sicherstellung einer robusten Infrastruktur, die modernen DevOps-Praktiken entspricht.

\bulletsend

Auch wenn meine aktuellen Frontend-Projekte primär auf React und Next.js basieren, handelt es sich bei beiden um TypeScript-basierte Component-Frameworks. Ich bin mit Angular, RxJS und Reactive Forms bestens vertraut und bringe große Motivation mit, diese Technologien schnell in meinen Arbeitsalltag zu integrieren. Meine umfangreiche Backend-Erfahrung mit NestJS und Node.js entspricht direkt Ihrem Tech-Stack. Ich arbeite gerne in cross-funktionalen Agile-Teams, in denen ich eng mit UX-Designerinnen, QA-Ingenieuren und Product Ownern zusammenarbeite, um Figma-Designs effizient in produktionsreife Software zu überführen.

Ich stehe ab sofort für eine Vollzeitstelle zur Verfügung und bin auch an Teilzeitmodellen interessiert. Ich befinde mich in Deutschland und bin innerhalb von zwei bis drei Wochen umzugsbereit. Zudem erfülle ich die Voraussetzungen für die Blaue Karte EU, sodass lediglich die Ausstellung eines regulären Arbeitsvertrags erforderlich ist, ohne zusätzlichen Visa-Sponsoring-Aufwand.

\closing{Mit freundlichen Grüßen,}{Radheem Bin Razi}

\end{document}
